\section{Datos y diseño de la evaluación}
\label{chap:diseno}

\subsection{Datos y tareas}
\label{sec:datos_tareas}

La evaluación utiliza dos conjuntos de datos con características y criterios de etiquetado diferentes, recogidos en la Tabla~\ref{tab:datasets_finales}. El simulador QRS/ST proporciona 100 latidos sintéticos con etiquetas que indican la presencia de tres rasgos morfológicos. Para la clasificación binaria, se considera positivo un caso cuando presenta simultáneamente RBBB, elevación del segmento ST e inversión de la onda T. Por su parte, los registros del HUCA utilizan una etiqueta clínica que indica la presencia o ausencia de síndrome de Brugada, sin identificar individualmente los rasgos del trazado. Aunque estas alteraciones pueden identificarse visualmente, su clasificación a partir de un único latido no equivale al diagnóstico de la enfermedad, que requiere una evaluación clínica más amplia.

\begin{table}[H]
  \centering
  \small
  \begin{tabularx}{\textwidth}{p{3cm}rrrr>{\raggedright\arraybackslash}X}
    \toprule
    Conjunto & Casos & Imágenes & Positivos & Negativos & Etiqueta \\
    \midrule
    Simulador QRS/ST & 100 & 100 & 20 & 80 & Presencia de los tres rasgos \\
    Datos HUCA & 317 & 317 & 116 & 201 & Etiqueta clínica \\
    \bottomrule
  \end{tabularx}
  \caption{Conjuntos de datos utilizados en la evaluación. Cada caso aporta una imagen correspondiente a un latido de V1.}
  \label{tab:datasets_finales}
  \label{tab:huca_datos_finales}
\end{table}

Los registros clínicos proceden del conjunto Brugada-HUCA \cite{costacortez2026brugada}, compuesto inicialmente por 363 pacientes. Siguiendo las recomendaciones de los médicos colaboradores, se excluyeron 46 de los 363 casos iniciales cuya clasificación clínica no estaba suficientemente definida en el momento de realizar el estudio. La muestra final quedó formada por 317 pacientes. De cada registro se seleccionó la derivación V1 y se extrajo un latido de la zona central, utilizando una ventana comprendida entre 300~ms antes y 600~ms después del pico R. Los detalles del procesamiento se recogen en el apéndice~\ref{app:configuracion}.

La Figura~\ref{fig:synthetic_real_comparison} muestra un ejemplo de cada conjunto para ilustrar las diferencias en la morfología y la presentación de los trazados. Las imágenes se han seleccionado de manera independiente y no corresponden a casos equivalentes, por lo que su comparación tiene únicamente una finalidad visual.

\begin{figure}[H]
  \centering
  \safeincludegraphics[width=0.92\textwidth]{results/fig_synthetic_real_comparison.pdf}
  \caption{Ejemplos de latidos del simulador y del HUCA en la derivación V1, seleccionados de forma independiente para ilustrar las diferencias entre ambos conjuntos.}
  \label{fig:synthetic_real_comparison}
\end{figure}

\subsubsection{Características de la muestra del HUCA}
\label{sec:procedencia_huca}

La muestra clínica contiene 116 pacientes etiquetados como positivos y 201 como negativos. Esta distribución presenta un desequilibrio entre clases que debe tenerse en cuenta durante la evaluación. Un clasificador que predijera siempre la clase mayoritaria alcanzaría una exactitud del 63,4~\%, pero no identificaría ningún caso positivo. Por ello, resulta necesario utilizar métricas que permitan valorar el rendimiento en ambas clases.

Aunque el procesamiento permite obtener imágenes de las derivaciones V1, V2 y V3, la evaluación se limita a V1, utilizando un único latido por paciente. Esta elección se fundamenta en la relevancia de las derivaciones precordiales derechas para la identificación de patrones asociados al síndrome de Brugada, especialmente las alteraciones del segmento ST y de la onda T \cite{zeppenfeld2022}. No obstante, el análisis de una sola derivación limita la información disponible respecto a una evaluación clínica completa.

\section{Implementación de los modelos}
\label{chap:modelos}

Los experimentos comparan dos enfoques de clasificación que utilizan imágenes de un latido de V1, pero aprovechan los ejemplos etiquetados de manera diferente. ResNet18 utiliza estos ejemplos durante el entrenamiento supervisado para ajustar sus pesos, mientras que Gemma~4 los incorpora como demostraciones en la consulta mediante aprendizaje en contexto, manteniendo sus parámetros fijos. La Figura~\ref{fig:diseno_experimental_resumen} representa ambos procedimientos y el apéndice~\ref{app:configuracion} recoge los detalles de su configuración.

\begin{figure}[H]
  \centering
  % Fuente editable: assets/results/fig_pipeline_experimental.drawio.
  \safeincludegraphics[width=0.96\textwidth]{results/fig_pipeline_experimental.pdf}
  \caption{Comparación de los procedimientos de clasificación utilizados. ResNet18 ajusta sus parámetros mediante entrenamiento supervisado, mientras que Gemma~4 utiliza demostraciones etiquetadas durante la inferencia sin modificar sus pesos.}
  \label{fig:diseno_experimental_resumen}
\end{figure}

\subsection{Gemma con aprendizaje en contexto}
\label{chap:vlm}

\subsubsection{Modelo y servicio}

Para los experimentos se utiliza Gemma~4 E4B, un modelo de visión y lenguaje de pesos abiertos desarrollado por Google. La variante \texttt{google/gemma-4-E4B-it} admite entradas de imagen y texto y está adaptada para seguir instrucciones \cite{gemma4card}. Durante la evaluación, los pesos del modelo permanecen fijos y se modifica el número y la procedencia de los ejemplos etiquetados incluidos en cada consulta.

La inferencia se ejecuta mediante vLLM~\cite{kwon2023vllm} en un clúster privado equipado con dos unidades de procesamiento gráfico (GPU) NVIDIA A100 y 96~GB de memoria de vídeo (VRAM). La temperatura se fija en cero para reducir la variabilidad de las respuestas. Los parámetros de inferencia se recogen en la Tabla~\ref{tab:config_vlm_appendix}.

\subsubsection{Prompt y respuesta}

Cada consulta comienza con una instrucción de sistema que define la tarea de clasificación y el formato de respuesta esperado. A continuación, se incorporan las demostraciones mediante pares de turnos. En cada uno, el mensaje del usuario contiene una imagen de ECG de la derivación V1 y el del asistente proporciona su etiqueta conocida en formato JSON. Finalmente, se añade la imagen que debe clasificarse, cuya etiqueta se mantiene oculta durante la inferencia y se utiliza posteriormente para evaluar la predicción.

La Figura~\ref{fig:icl_context_window} ilustra este procedimiento para la clasificación de los tres rasgos morfológicos del simulador. La selección de demostraciones sigue las condiciones estándar y balanceada descritas en la sección~\ref{sec:condiciones_vlm}.

\begin{figure}[H]
  \centering
  % Fuente editable: assets/results/vlm_icl_context_window.drawio.
  \safeincludegraphics[width=\textwidth]{results/vlm_icl_context_window.pdf}
  \caption{Estructura de una consulta de aprendizaje en contexto con demostraciones etiquetadas y una imagen de prueba para la clasificación de los tres rasgos morfológicos del simulador.}
  \label{fig:icl_context_window}
\end{figure}

El contenido de la respuesta depende de la tarea evaluada. En el simulador, Gemma~4 debe identificar por separado la presencia de RBBB, elevación del segmento ST e inversión de la onda T. A partir de estas tres etiquetas se obtiene también una clasificación binaria, considerando positivo un caso cuando presenta los tres rasgos simultáneamente. En los experimentos con registros del HUCA, el modelo genera directamente una etiqueta positiva o negativa correspondiente a la clasificación clínica de Brugada.

El apéndice~\ref{app:prompts} recoge las instrucciones y los formatos de salida empleados en cada experimento.

\subsection{CNN ResNet18}
\label{chap:cnn_enfoque}

\subsubsection{Arquitectura y entrenamiento}
\label{sec:cnn_entrenamiento}

ResNet18 parte de los pesos preentrenados en ImageNet-1K que proporciona Torchvision~\cite{torchvisionresnet18}. Sus bloques convolucionales extraen características de las imágenes que la última capa utiliza para generar las predicciones. En los experimentos con el simulador, esta capa tiene tres unidades de salida, una por cada rasgo morfológico, mientras que en los del HUCA se utiliza una única unidad para la etiqueta clínica. En ambos casos se ajustan los pesos de toda la red. La Figura~\ref{fig:cnn_convolutional_pipeline} muestra el procedimiento de entrenamiento y evaluación para el simulador.

\begin{figure}[H]
  \centering
  % Fuente editable: assets/results/cnn_convolutional_pipeline.drawio.
  \safeincludegraphics[width=0.96\textwidth]{results/cnn_convolutional_pipeline.pdf}
  \caption{Entrenamiento supervisado y evaluación de ResNet18 para la clasificación de los tres rasgos morfológicos del simulador.}
  \label{fig:cnn_convolutional_pipeline}
\end{figure}

El preprocesamiento convierte las imágenes a escala de grises, las redimensiona a 224~×~224 píxeles y replica la intensidad en tres canales, como requiere ResNet18. La media y la desviación estándar se calculan exclusivamente con los datos de entrenamiento y se utilizan también para normalizar las imágenes de validación y prueba.

\begin{samepage}
La función de pérdida utilizada es la entropía cruzada binaria, aplicada a cada uno de los tres rasgos del simulador o a la única etiqueta del HUCA. Para compensar el desequilibrio entre clases, los errores de los ejemplos positivos se ponderan mediante el cociente entre negativos y positivos del conjunto de entrenamiento. En la selección balanceada del HUCA, este peso es igual a uno.
\par
\end{samepage}

El entrenamiento utiliza el optimizador AdamW \cite{loshchilov2019}, que actualiza los pesos de la red a partir de los gradientes calculados en cada lote de imágenes. La tasa de aprendizaje, el tamaño del lote y los demás parámetros se recogen en la Tabla~\ref{tab:config_cnn_appendix}.

\subsubsection{Predicción y selección del umbral}
\label{sec:umbral_cnn}

Cada salida de ResNet18 es una puntuación denominada \textit{logit}, que se transforma en una probabilidad estimada mediante la función sigmoide. En el simulador se obtiene una probabilidad para cada rasgo morfológico, mientras que en el HUCA se calcula una única probabilidad asociada a Brugada. Ambas operaciones se expresan como
\begin{equation}
  \begin{gathered}
    p_j = \frac{1}{1+\exp(-z_j)},\quad j=1,2,3\\
    \text{Simulador}
  \end{gathered}
  \qquad
  \begin{gathered}
    p_B = \frac{1}{1+\exp(-z_B)}\\
    \text{Datos clínicos}
  \end{gathered}
  \label{eq:probabilidades_cnn}
\end{equation}
Los tres rasgos del simulador pueden coexistir, por lo que sus probabilidades no tienen que sumar uno. La clasificación binaria derivada considera positivo un caso cuando las tres probabilidades alcanzan el umbral establecido. En los registros del HUCA, la decisión se obtiene directamente a partir de \(p_B\).

El umbral se selecciona buscando la mayor exactitud balanceada (BA) sobre los datos de validación. Para la clasificación conjunta del simulador se utiliza la menor de las tres probabilidades, ya que todas deben superar el mismo umbral.

\begingroup
\widowpenalty=10000
En los experimentos con datos clínicos, la selección del umbral también interviene durante el entrenamiento. Al finalizar cada época, se evalúa la red con los ocho pacientes de validación definidos en la sección~\ref{sec:validacion_comun} y se calcula la mejor BA obtenida. Este resultado permite seleccionar los pesos sin utilizar al paciente reservado para prueba.

El entrenamiento se interrumpe mediante \textit{early stopping} cuando transcurren cinco épocas sin mejora, con un máximo de veinte. Una vez recuperados los mejores pesos, se fija el umbral con los mismos pacientes de validación y se realiza la predicción del caso de prueba. Los resultados finales se obtienen agregando las predicciones de todos los pliegues.
\par
\endgroup

El ajuste de la CNN se realiza con pocos ejemplos, por lo que puede existir riesgo de sobreajuste.

\subsubsection{Adaptación de dominio}
\label{sec:adaptacion_cnn}

Se evalúan cuatro métodos de adaptación de dominio para aprovechar imágenes clínicas del HUCA sin utilizar sus etiquetas, manteniendo la clasificación de los tres rasgos morfológicos del simulador.

SimCLR permite preparar el extractor de características antes del entrenamiento supervisado. A partir de dos versiones aumentadas de cada imagen clínica, aprende representaciones resistentes a variaciones de brillo, pequeños desplazamientos, ruido y desenfoque. Los pesos obtenidos se utilizan después para inicializar ResNet18, que se entrena con las etiquetas del simulador.

CORAL y MMD incorporan la adaptación durante el entrenamiento supervisado, combinando imágenes etiquetadas del simulador con imágenes clínicas sin etiquetar. CORAL aproxima las covarianzas de las características de ambos conjuntos, mientras que MMD reduce las diferencias entre sus distribuciones mediante kernels de base radial a varias escalas. La contribución de ambos métodos aumenta progresivamente durante el entrenamiento.

DANN utiliza también imágenes de ambos conjuntos, pero incorpora un clasificador auxiliar que intenta distinguir su procedencia. Mediante inversión de gradiente, el extractor aprende características que dificultan esta distinción mientras continúa clasificando los rasgos del simulador. El clasificador auxiliar se elimina durante la inferencia.

En las cuatro variantes se utilizan imágenes clínicas de todos los pacientes, incluidos los reservados posteriormente para prueba, sin acceder a sus etiquetas. Esta configuración corresponde a una adaptación de dominio transductiva y permite estudiar si aprovechar las características visuales del HUCA mejora la transferencia desde el simulador. La Tabla~\ref{tab:config_domain_appendix} recoge los parámetros y las condiciones de uso de los datos.

\subsubsection{Inspección con Grad-CAM}

Grad-CAM complementa las métricas de clasificación mediante mapas que muestran las regiones de la imagen asociadas a una predicción. En este trabajo se aplica al último bloque convolucional de ResNet18, antes de la reducción de las características para su clasificación. El mapa se calcula a partir de los gradientes de la primera salida y se redimensiona al tamaño de la imagen original. En la CNN entrenada con el simulador, esta salida corresponde a RBBB.

\begin{samepage}
La Figura~\ref{fig:gradcam_sintetico_modelos} muestra el trazado original, el mapa Grad-CAM normalizado y su superposición. En el ejemplo se observa una región amplia asociada a RBBB que comprende el complejo QRS y parte del trazado posterior. Esta inspección permite identificar las zonas relacionadas con la predicción, aunque no determina por sí sola qué características clínicas ha reconocido el modelo, como se explica en la sección~\ref{sec:explicabilidad}.
\par
\end{samepage}

\begin{figure}[H]
  \centering
  \safeincludegraphics[width=0.92\textwidth]{results/cnn_simulator_qrs_gradcam_example.pdf}
  \caption{Inspección de la salida RBBB de ResNet18 mediante Grad-CAM. Se muestran el trazado V1, el mapa normalizado y su superposición.}
  \label{fig:gradcam_sintetico_modelos}
\end{figure}



\section{Protocolo de evaluación}

\subsection{Escenarios de comparación}
\label{sec:escenarios}

La evaluación comprende tres escenarios que permiten estudiar el rendimiento de los modelos con datos sintéticos y clínicos, así como su capacidad de transferencia entre ambos conjuntos.

\begin{enumerate}
  \item \textbf{Aprendizaje con el simulador.} Gemma~4 recibe demostraciones etiquetadas con los tres rasgos morfológicos y clasifica nuevos casos sintéticos. ResNet18 utiliza esas etiquetas para su entrenamiento supervisado.
  \item \textbf{Transferencia del simulador al HUCA.} Gemma~4 recibe demostraciones del simulador y clasifica ECG clínicos. ResNet18 se entrena con datos sintéticos y se evalúa sobre registros del HUCA.
  \item \textbf{Aprendizaje con pacientes del HUCA.} Gemma~4 utiliza demostraciones con etiquetas clínicas del hospital, mientras que ResNet18 se entrena directamente con esas mismas etiquetas.
\end{enumerate}

\subsection{Selección de ejemplos y validación}
\label{sec:validacion_comun}

Se utiliza validación cruzada dejando un paciente fuera (LOOCV). En cada pliegue se reserva un paciente para prueba y se seleccionan los ejemplos entre los restantes, evitando que aparezca en el entrenamiento, la validación o las demostraciones de Gemma~4.

Se denomina \(K\) al número de demostraciones que recibe Gemma~4 o al número de casos de entrenamiento de ResNet18. Se evalúan \(K=\{2,4,8,16,32\}\) demostraciones para Gemma~4 y el mismo número de pacientes de entrenamiento para ResNet18. Cada configuración se repite con las semillas 42, 123 y 2026, que controlan la selección aleatoria de ejemplos y, en la CNN, la aleatoriedad del entrenamiento. Ambos modelos se evalúan sobre los mismos pacientes y seleccionan sus ejemplos del mismo conjunto disponible.

Las variantes de adaptación de dominio también utilizan imágenes clínicas sin etiquetar, incluidas aquellas que posteriormente se evalúan. Este procedimiento transductivo se describe en la sección~\ref{sec:adaptacion_cnn} y en la Tabla~\ref{tab:config_domain_appendix}.

En los experimentos con el HUCA, ResNet18 reserva además ocho pacientes etiquetados, cuatro por clase, para seleccionar el umbral y aplicar \textit{early stopping}. Estos pacientes son distintos de los utilizados para entrenamiento y prueba. Gemma~4 no requiere este conjunto adicional de validación.

\clearpage
\subsubsection{Secuencia de un pliegue clínico}

Para cada paciente reservado, Gemma~4 construye una consulta con las instrucciones, las demostraciones etiquetadas y la imagen de prueba. La etiqueta de este último caso permanece oculta hasta la evaluación.

ResNet18, por su parte, se inicializa con los pesos preentrenados y se ajusta utilizando los pacientes de entrenamiento del pliegue. Los casos de validación permiten seleccionar la época y el umbral antes de clasificar al paciente reservado. La Figura~\ref{fig:clinical_fold_protocol} resume ambos procedimientos.

\begin{figure}[H]
  \centering
  \safeincludegraphics[width=\textwidth]{results/clinical_fold_protocol.pdf}
  \caption{Procedimiento de evaluación de un pliegue con los datos clínicos del HUCA.}
  \label{fig:clinical_fold_protocol}
\end{figure}

La CNN utiliza \(K+8\) pacientes etiquetados por pliegue, frente a los \(K\) empleados como demostraciones por Gemma~4. Esta diferencia corresponde a los ocho pacientes adicionales de validación, cuyos datos no intervienen en la actualización de pesos. Considerando las selecciones estándar y balanceadas, los cinco valores de \(K\) y las tres semillas, la evaluación clínica de ResNet18 comprende 30 ejecuciones de 317 pliegues cada una.

\subsubsection{Distribución de los ejemplos por clase}
\label{sec:condiciones_vlm}
\label{sec:origen_contexto_vlm}

La proporción de ejemplos positivos y negativos puede influir en el rendimiento cuando se dispone de pocos casos. Para estudiar este efecto se comparan dos estrategias de selección. La \textbf{estándar} garantiza al menos un positivo y un negativo para \(K\geq2\) y completa el grupo de forma seudoaleatoria. La \textbf{balanceada} selecciona el mismo número de ejemplos de cada clase y, si no hay suficientes de alguna, completa el grupo con casos de la otra.

\begin{samepage}
Ambas estrategias se aplican a las demostraciones de Gemma~4 y al entrenamiento de ResNet18 con el HUCA. En la CNN, la selección estándar conserva los pacientes definidos en el plan de pliegues, mientras que la balanceada mantiene los ocho pacientes de validación y selecciona nuevamente los de entrenamiento. Gemma~4 elige sus demostraciones entre los pacientes disponibles, aplicando la distribución correspondiente.
\par
\end{samepage}

También se evalúa Gemma~4 en condiciones \emph{zero-shot} (\(K=0\)), utilizando únicamente las instrucciones y la imagen de prueba, sin demostraciones. Esta configuración permite comparar el rendimiento sin ejemplos con el obtenido mediante las selecciones estándar y balanceada. Para \(K>0\), la procedencia de los ejemplos depende del escenario experimental.

\subsection{Métricas y agregación}
\label{sec:metricas}

La exactitud balanceada (BA) se utiliza como métrica principal, ya que asigna el mismo peso a la sensibilidad y la especificidad, independientemente de la proporción de cada clase. Se calcula mediante la expresión
\[
  \mathrm{BA}=\frac{1}{2}\left(\frac{TP}{TP+FN}+\frac{TN}{TN+FP}\right).
\]

La sensibilidad mide la proporción de positivos correctamente identificados y la especificidad, la de negativos. Los verdaderos positivos (TP) y verdaderos negativos (TN) representan los aciertos, mientras que los falsos negativos (FN) y falsos positivos (FP) corresponden a los errores de clasificación.

La métrica F1 combina precisión y sensibilidad mediante su media armónica. Para las predicciones probabilísticas se calcula también el área bajo la curva ROC (ROC-AUC), que mide la capacidad de discriminación entre clases al variar el umbral. La precisión media (\textit{average precision}) pondera la precisión en cada umbral por el incremento de sensibilidad respecto al anterior~\cite{sklearnaverageprecision}. Las matrices de confusión permiten examinar la distribución de aciertos y errores de cada modelo.

En una posible aplicación de triaje, sensibilidad y especificidad deben interpretarse conjuntamente, ya que aumentar la detección de casos positivos puede incrementar las falsas alertas y la carga de revisión de los especialistas.
